\documentclass[11pt]{article}
\usepackage{xltxtra}
\usepackage{bookmark}
\usepackage{hyperref}
\hypersetup{hidelinks}
\usepackage{url}
\urlstyle{tt}
\usepackage{multicol}
\usepackage{xcolor}
\usepackage{calc}
\usepackage{graphicx}
\usepackage{tikz}
\usetikzlibrary{calc}
\usepackage{fontspec}
\usepackage{xeCJK}
\usepackage{relsize}
\usepackage{xspace}
\usepackage{fontawesome5}
\usepackage{titlesec}
\usepackage{enumitem}
\usepackage{siunitx}
\usepackage{amssymb}
\usepackage{tabularx}
\usepackage{multicol}
\usepackage{fontspec}
\usepackage{fancybox}
\usepackage{float}

% 一些小设置，参考自https://github.com/LeyuDame/BNUCV/tree/main
% 取消中文字符与数字之间的间隔
\CJKsetecglue{}
\protected\def\Cpp{{C\nolinebreak[4]\hspace{-.05em}\raisebox{.28ex}{\relsize{-1}++}}\xspace}	% 这是个更好看的C++写法，你直接写C++的话，+号会很大，可以使用\Cpp来代替
\setlength{\parindent}{0pt}							% 取消全局段落缩进
\pagenumbering{gobble}								% 取消页码显示
%\setlist{noitemsep}									% 禁用列表中项目之间的额外垂直间距，但保留列表周围的间距
%\setlist{nosep}										% 禁用列表中项目之间的额外垂直间距及列表周围的间距
\setlist[itemize]{topsep=0em, leftmargin=*}		% 增加了itemize顶部间距
\setlist[enumerate]{topsep=0em, leftmargin=*}	% 增加了enumerate顶部间距

\titleformat{\section}					    % 将原标题前面的数字取消了
  {\LARGE\bfseries\raggedright} 		      % 字体改为LARGE，bold，左对齐
  {}{0em}                      			  % 可用于添加全局标题前缀
  {}                           			  % 可用于添加代码
  [{\color{WHU_Blue}\titlerule}]            % 标题下方加一条线
\titlespacing*{\section}{0cm}{*1.2}{*1.2}	% 标题左边留白，上方1.2倍，下方1.2倍

\titleformat{\subsection}				    % 将原二级标题前面的数字取消了
  {\large\bfseries\raggedright} 		      % 字体改为large，bold，左对齐
  {}{0em}                      			  % 可用于添加全局二级标题前缀
  {}                           			  % 可用于添加代码
  []
\titlespacing*{\subsection}{0cm}{*1.2}{*1.2}% 二级标题左边留白，上方1.2倍，下方1.2倍
% 页面大小与页边距，按需求调整
\usepackage[
	a4paper,
	left=1.2cm,
	right=1.2cm,
	top=1.5cm,
	bottom=1cm,
	nohead
]{geometry}
% 这里把表格的行间距调成1.2倍了
\renewcommand{\arraystretch}{1.2}
% 这里把正文的行间距调成1.2倍了
\linespread{1.2}
% 中文字符间距
\renewcommand{\CJKglue}{\hskip 0.05em}

% 自定义宋体（狮尾四季春），详见readme
% 英文字体
\setmainfont[
    Path=fonts/,
    Extension=.ttf,
    BoldFont=* Bold,
]{SweiSpring}
% 中文字体
\setCJKmainfont[
    Path=fonts/,
    Extension=.ttf,
    BoldFont=* Bold,
]{SweiSpring}

% 自定义主题色
% 武大蓝（来自官方配色方案，参考https://www.whu.edu.cn/xxgk/wdbs.htm）
\definecolor{WHU_Blue}{RGB}{0, 37, 84}
% 西大蓝 （来自官方VIS配色方案，参考http://dwxcbwlgzb.swu.edu.cn/info/1017/1016.htm）
\definecolor{SWU_Blue}{HTML}{353F6C}



%%%%%%%%%%%%%%%%%%%%%%%%%%%%%%%
%%%%%%%%%  页眉页脚  %%%%%%%%%%%
%%%%%%%%%%%%%%%%%%%%%%%%%%%%%%%
% 定义填写学院信息
\newcommand{\school}{人工智能学院 | College of AI} 

% 定义联系方式
\newcommand{\contact}{
    % 根据个人喜好选择字号
    % \small              % 小
    % \footnotesize       % 更小
    \scriptsize         % 再小一号
    \textcolor{white}{
        % 邮箱 \href中写动态链接 后面写为文本，支持点击文本自动跳转
        \faEnvelope \quad \href{mailto:warmthwind@163.com}{warmthwind@163.com}
        \hspace{4em}
        % GitHub
        \faGithub \quad \href{https://github.com/MrKeanu-V}{MrKeanu-V}
        \hspace{4em}
        % 手机号
        \faPhone \quad  187****4104（微信同号）
        % 添加其他内容如微信、GitHub等格式如下，使用\fa命令添加不同icon，细节请搜索Font Awesome图标库
        % \hspace{4em}
        % \faGithub \quad \href{https://github.com/xxxx}{https://github.com/xxxx}
    }
}



%%%%%%%%%%%%%%%%%%%%%%%%%%%%%%%
%%%%%%%%%  简历正文  %%%%%%%%%%%
% 写简历还有项目细节时不要泛泛而谈太笼统，要应用STAR原则，即Situation（情景）、Task（任务）、Action（行动）和Result（结果）四个英文单词的首字母组合。
% S指的是situation，事情是在什么情况下发生
% T指的是task，你是如何明确你的目标的
% A指的是action，针对这样的情况分析，你采用了什么行动方式
% R指的是result，结果怎样，在这样的情况下你学习到了什么
%%%%%%%%%%%%%%%%%%%%%%%%%%%%%%%
\begin{document}
% 如果有多页简历，请把页眉页脚和背景复制粘贴到第二页的内容之前
   %%%%%%%%%  1、页眉页脚  %%%%%%%%%%%
    % 页眉：校标组合+学院名
    \begin{tikzpicture}[remember picture, overlay]
        \node[anchor=north, inner sep=0pt](header) at (current page.north){
            \includegraphics[width=\paperwidth]{images/header.png}
        };
        \node[anchor=west](school_logo) at (header.west){
            \hspace{0.5cm}
            \includegraphics[width=0.15\textwidth]{images/swu_logo_1_AIenhanced.png}
        };
        \node[anchor=east](school_name) at(header.east){
            \textcolor{white}{\textbf{\school}}
            \hspace{0.5cm}
        };
    \end{tikzpicture}
    \vspace{-3.5em}

    % 页脚：为文首声明的\contact，
    \begin{tikzpicture}[remember picture, overlay]
        \node[anchor=south, inner sep=0pt](footer) at (current page.south){
            \includegraphics[width=\paperwidth]{images/footer.png}
        };
        % 需要更改时直接更改文首的页眉页脚即可
        \node[anchor=center] at(footer.center){\contact};
    \end{tikzpicture}

    % 背景：校徽背景
    \begin{tikzpicture}[remember picture, overlay]
        \node[opacity=0.05] at(current page.center){
            \includegraphics[width=0.7\paperwidth, keepaspectratio]{images/swu_logo_big_blue.png}
        };
    \end{tikzpicture}


    %%%%%%%%%  2、简历正文  %%%%%%%%%%%
    % 个人信息
    \begin{figure}[h]
        % 左半边，信息，比例占行宽82%，可以自己调，非必要不建议调整
        \begin{minipage}{0.82\textwidth}
            \section{\makebox[\widthof{\faAddressCard}][cc]{\color{SWU_Blue}{\faAddressCard}}\quad 个人信息}
            \begin{tabularx}{\linewidth}{p{\widthof{出生日期:}}Xp{\widthof{政治面貌:}}X}
                姓\ \ \ \ \ \ \ \ \ 名: & 方鸿渐 & 
                性\ \ \ \ \ \ \ \ \ 别: & 男  \\
                出生年月: & XXXX.XX.XX & 
                籍\ \ \ \ \ \ \ \ \ 贯: & 重庆  \\
                %% 想多加几行的话，就按上面的格式自行补充
                %% 想加粗的话\textbf{}
                %% 想多加几列的话，把\begin{tabularx}{\textwidth}{这里}的内容改一下，可以自己搜一下tabularx怎么用，也可以问gpt/文心一言/讯飞。
            \end{tabularx}
        \end{minipage}
    \hspace{2em}
    % 右半边，照片，比例占行宽12%，可以自己调
    % images/avatar.png 替换成你证件照的路径。
    \begin{minipage}{0.12\textwidth}
        \setlength{\fboxsep}{0pt}
        \doublebox{\includegraphics[width=\linewidth]{images/my_avatar.png}}
    \end{minipage}
    \end{figure}
    \vspace{-1.5em}


    % 教育背景（本科生）
    % 这类\fa开头的都是font awesome里的logo，想换成其他logo的话，可以搜一下fontawsome支持的icon指令。
    \section{\makebox[\widthof{\faGraduationCap}][c]{\color{SWU_Blue}{\faGraduationCap}}\quad 教育背景}

    % 硕士研究生
    {\large \textbf{西南大学}}  学士 \hfill {重庆} \\
    {{人工智能学院}}·人工智能 \hfill {20XX年9月-20XX年6月} \\
    {主修课程}：课程1、课程2、课程3、课程4\ 等
    \hfill {GPA}：4.0/4.0

    % 博士研究生
    \vspace{0.5em}
    {\large \textbf{西南大学}}  硕士 \hfill {重庆} \\
    {{人工智能学院}}·人工智能 \hfill {20XX年9月-20XX年6月} \\
    \textbf{研究方向}：方向1、方向2、方向3、方向4\ 等
    \hfill {GPA}：100/100

    \vspace{0.5em}
    {\large \textbf{西南大学}}  博士 \hfill {重庆} \\
    {{人工智能学院}}，导师：\href{导师的个人主页.site}{导师名字}\ 导师职称 \hfill {20XX年9月-至今} \\
    \textbf{研究方向}：方向1、方向2、方向3、方向4\ 等


    % 科研著作（研究生）
    \section{\makebox[\widthof{\faGraduationCap}][c]{\color{SWU_Blue}{\faFlask}}\quad 科研成果}
    
    \vspace{0.5em}
    This is Another Paper. \\
    \textbf{Mingzi Nide}, Shidi Nide, Daoshi Nide. \hfill 
    在投 \textbf{Conference B} （CCF-B类会议）

    \vspace{0.5em}
    This is A Journal Paper.\\
    \textbf{Mingzi Nide}, Shixiong Nide, Daoshi Nide. \hfill 
    在投 \textbf{Journal C} （SCI-2区）


    % 项目经历\科研经历\项目与教学（标题请根据需要修改）
    \section{\makebox[\widthof{\faChalkboardTeacher}][c]{\color{SWU_Blue}{\faChalkboardTeacher}}\quad 横向项目}
    % 项目模板1
    \vspace{0.5em}
    {\large{\textbf{项目名称}}} \hfill {横向/纵向项目-已完结/进行中}\\
    \textbf{你在项目中扮演的角色} \hfill 2020年9月至2021年9月\\
    项目简介：XXX，XXX。
    
    % 项目模板2
    \vspace{0.5em}
    {\large{\textbf{国企单位-基于视觉识别算法改进的厂区安全检测平台}}} 项目负责人 \hfill {2023年10月至2024年1月}\\
    主要内容：数据集制作;模型复现训练;核心算法实现;GUI开发;专利撰写\ 等

    % 技能特长（标题根据个人需求修改）
    \section{\makebox[\widthof{\faWrench}][c]{\color{SWU_Blue}{\faWrench}}\quad 技能荣誉 }
    \vspace{0.5em}
    \begin{itemize}
        \item 英语：六级800分、托福200分
        \item 编程：Python, R, MATLAB, C
        \item 证书：CET-6, 教师资格证书
        \item 其他： 专利X项, 软著X项,其他奖项若干
    \end{itemize}
    

    %%%%%%%%%  3、新加一页的模板，根据需要可添加  %%%%%%%%%%%
    % \newpage

    % 页眉页脚不要删。
    % % 页眉：校标组合+学院名
    % \begin{tikzpicture}[remember picture, overlay]
    %     \node[anchor=north, inner sep=0pt](header) at (current page.north){
    %         \includegraphics[width=\paperwidth]{images/header.png}
    %     };
    %     \node[anchor=west](school_logo) at (header.west){
    %         \hspace{0.5cm}
    %         \includegraphics[width=0.15\textwidth]{images/swu_logo_1_AIenhanced.png}
    %     };
    %     \node[anchor=east](school_name) at(header.east){
    %         \textcolor{white}{\textbf{\school}}
    %         \hspace{0.5cm}
    %     };
    % \end{tikzpicture}
    % \vspace{-4em}
    
    % % 页脚，联系方式
    % \begin{tikzpicture}[remember picture, overlay]
    %     \node[anchor = south, inner sep=0pt] at (current page.south){
    %         \includegraphics[width=\paperwidth]{images/footer.png}
    %     };
    %     % 联系方式
    %     \node[anchor=center] at(footer.center){\contact};
    % \end{tikzpicture}
    
    % % 背景
    % \begin{tikzpicture}[remember picture, overlay]
    %     \node[opacity=0.1] at(current page.center){
    %         \includegraphics[width=0.7\paperwidth, keepaspectratio]{images/swu_logo_big_blue.png}
    %     };
    % \end{tikzpicture}


    % % 竞赛经历
    % \section{\makebox[\widthof{\faTrophy}][c]{\color{SWU_Blue}{\faTrophy}}\quad 竞赛经历}
    % \vspace{-1em}
    % \begin{table}[h!]
    %     \begin{tabularx}{\textwidth}{Xp{\widthof{第零负责人}}p{\widthof{国家级-第100名}}p{\widthof{2030年13月}}}
    %         \textbf{比赛1} & 第一负责人 & 国家级-第10名 & 2023年4月 \\
    %         \textbf{比赛2} & 个人参赛 & 国家级-一等奖 & 2023年8月\\
    %         \textbf{比赛3} & 个人参赛 & 省级-一等奖 & 2022年12月\\
    %         % 同理，可以自己加
    %     \end{tabularx}
    % \end{table}

    % % 技能特长
    % \section{\makebox[\widthof{\faWrench}][c]{\color{SWU_Blue}{\faWrench}}\quad 技能特长}
    % \vspace{0.5em}
    % \begin{itemize}
    %     \item 熟练使用\Cpp 、Python、Matlab编程语言。
    %     \item 熟悉Windows与Linux端开发。
    %     \item 熟练使用Tensorflow，Pytorch等深度学习框架。
    %     \item 熟练掌握\Cpp 与Python环境下OpenCV与Qt应用的开发，且熟练使用Qt Creator软件。
    %     \item 熟练使用Altium Designer与LCEDA进行封装绘制与板子设计。
    %     \item 熟练使用Keil，Arduino IDE等集成开发软件。
    %     \item 了解模式识别，强化学习，遗传算法，知识蒸馏等相关概念。
    % \end{itemize}

    % % 所获荣誉
    % \section{\makebox[\widthof{\faStar}][c]{\color{SWU_Blue}{\faStar}}\quad 所获荣誉}
    % \vspace{-1em}
    % \begin{multicols}{2}
    %     \begin{itemize}
    %         \item 某年学业先进个人
    %         \item 某年某奖学金某等奖
    %         \item 某大使
    %         \item 某年某奖学金某等奖
    %         \item 某年优秀团员称号
    %         \item 某年某称号
    %     \end{itemize}
    % \end{multicols}

    % % 其他
    % \section{\makebox[\widthof{\faInfo}][c]{\color{SWU_Blue}{\faInfo}}\quad 其他}
    % \begin{itemize}
    %     \item 英语水平-CET6级xxx分
    %     \item 计算机几级证书
    %     \item xx几级证书
    %     \item 技术博客: 某网址
    %     \item 教师资格证：xxx
    %     \item 普通话证书：几级几等
    %     \item 文字排版：\LaTeX
    % \end{itemize}

\end{document}
